\subsection{Modelos de ensamble}
Se compararon Random Forest (RF) e HistGradientBoosting (HGB) utilizando el mismo objetivo, predictores y particiones cronológicas. La etiqueta indica si la tasa mediana del bloque siguiente supera el percentil 75 de las 1008 tasas observadas hasta el bloque actual. Después del filtro desde enero de 2016 y de eliminar las filas sin historia completa, se obtuvieron 418\,720 ejemplos: 293\,104 de entrenamiento, 62\,808 de validación y 62\,808 de prueba reservada.

Los 33 predictores incluyen estadísticas del bloque actual, rezagos de comisiones de 1, 6 y 144 bloques, medias y desviaciones de 6, 144 y 1008 bloques, el umbral local y razones relativas de comisión, indicadores de actividad y variables cíclicas de calendario. Todas las ventanas terminan en el bloque actual. Los intervalos temporales negativos se marcaron explícitamente y se trataron como faltantes. La imputación por mediana se ajustó dentro de cada conjunto de entrenamiento. No se estandarizaron las entradas de los árboles ni se eliminaron automáticamente valores extremos.

RF combina 100 árboles entrenados con muestras bootstrap de tamaño equivalente a la mitad del conjunto de ajuste y considera la raíz cuadrada del número de predictores en cada división. Se exploraron tres combinaciones de profundidad máxima y mínimo de observaciones por hoja. HGB construye árboles secuenciales sobre representaciones discretizadas de los predictores; se compararon tres configuraciones de número de iteraciones, hojas, tasa de aprendizaje, tamaño mínimo de hoja y regularización L2. La búsqueda acotada y sus configuraciones completas se incluyen en el repositorio. Se utilizó semilla 42 y se desactivó el criterio de parada temprana interno de HGB para evitar una separación aleatoria adicional.

\subsection{Diseño experimental}
La selección utilizó tres particiones temporales expansivas dentro del entrenamiento, con 40\,000 ejemplos de validación interna y un margen conservador de 1008 bloques en cada partición, siguiendo el diseño propuesto en P1. Se seleccionó por familia el candidato con mayor F1 medio usando un umbral inicial de probabilidad de 0.5. Después se ajustó el umbral de decisión con sus probabilidades fuera de fold, en una malla de 0.10 a 0.90 con paso 0.025. Estas puntuaciones intervienen en la selección y no se interpretan como una evaluación independiente.

Los modelos elegidos se reajustaron sobre 292\,096 ejemplos, reservando las últimas 1008 filas del entrenamiento nominal como margen frente a validación. Ambos se evaluaron sobre las mismas 62\,808 observaciones posteriores. Se reportaron F1 de la clase alta, precisión, recall, average precision (AP), ROC-AUC y matrices de confusión. AP se distingue del área trapezoidal bajo la curva precisión--recall, también disponible en los archivos de resultados. La evaluación de prueba se mantiene pendiente hasta congelar el protocolo conjunto.

\subsection{Diagnóstico e interpretación}
Se comparó el F1 en ajuste y validación como diagnóstico de generalización, reconociendo que la diferencia también puede responder a cambios temporales. Se obtuvo importancia por reducción de impureza para RF, siguiendo el ejemplo del curso. Esta medida es descriptiva y puede verse afectada por cardinalidad y correlación entre variables; no implica causalidad. HGB se interpreta a través de las métricas, curvas y matrices de confusión. El estudio de ablación con reentrenamiento corresponde al integrante 3.
